% ==========================================================================
%  Appendix B — Benchmark Data
% ==========================================================================
\chapter{Benchmark Data}
\label{app:benchmarks}

This appendix presents the complete benchmark data and methodology
underlying the performance claims in Chapters~\ref{ch:v2redesign}
and~\ref{ch:conclusion}.  All measurements were taken on the
\texttt{mjbommar/linux} tree (branch \texttt{next}), kernel commit
\texttt{9ccbf5e7bb58}, with the host CPU governor set to
\texttt{performance}.  Hosts span Skylake-S through Zen~4 and
Alder Lake.


% ------------------------------------------------------------------
\section{Methodology}
\label{sec:bench-method}
% ------------------------------------------------------------------

The benchmark suite uses three tiers, each measuring a different
aspect of backend performance:

\begin{description}[style=nextline,leftmargin=2em]
  \item[Micro tier.]
    A tight \texttt{getpid()} loop bracketed by \texttt{RDTSC}.
    Measures pure trap-mechanism cost: the round-trip from guest
    userspace through the backend to the \UML{} kernel and back.
    This tier is dominated by the syscall interception mechanism and
    is insensitive to kernel processing time.

  \item[Application tier (Python startup).]
    A canned Python interpreter startup with a mixed-syscall workload
    (\texttt{python3 -c "import hashlib"}).  Exercises
    \texttt{mmap}, \texttt{mprotect}, \texttt{open}, \texttt{read},
    \texttt{fstat}, \texttt{brk}, process creation, dynamic linking,
    and C-extension loading.

  \item[Stress tier (mt-mini).]
    An SMP stress workload with T=8 threads on ncpus=4, performing
    concurrent \texttt{mmap}/\texttt{munmap}/\texttt{mprotect}
    operations with strict memory verification after each iteration.
    Measures both backend performance and SMP correctness under
    concurrency pressure.
\end{description}

\noindent
The benchmark unit is speedup over seccomp on the same host.  This
ratio cancels most CPU-frequency and microarchitecture noise and
ties the claim to a backend decision rather than a hardware selection.


% ------------------------------------------------------------------
\section{Per-Syscall Latency}
\label{sec:bench-syscall}
% ------------------------------------------------------------------

Table~\ref{tab:syscall-latency} summarizes the round-trip cost of a
single \syscall{getpid} invocation under each backend mechanism,
contextualized against a native Linux host syscall baseline and
the LKL direct-call reference.

\begin{table}[htbp]
\centering
\caption{Per-syscall (\texttt{getpid}) round-trip latency by backend.}
\label{tab:syscall-latency}
\begin{tabular}{@{}llrrl@{}}
\toprule
Backend & Mode & Cycles & Approx.\ ns & Notes \\
\midrule
Native host syscall & ring 0 entry & $\sim$150 & $\sim$50 & Baseline (not UML) \\
\KVMvTwo{} + gadget & In-guest LSTAR & $\sim$90 & $\sim$25 & No VMEXIT; 10 trivial NRs \\
\KVMvTwo{} (exit path) & \texttt{KVM\_EXIT\_IO} & $\sim$36,800 & $\sim$150 & IO-port exit path \\
\Seccomp{} & SIGSYS + futex & $\sim$110,000 & 300--500 & Default backend (6.16+) \\
ptrace (archived) & \texttt{PTRACE\_SYSEMU} & 3,000--15,000 & 1,000--5,000 & Removed in memo 25 R11 \\
\midrule
LKL (direct call) & function call & $\sim$18 & $\sim$6 & No ring split; not comparable \\
\bottomrule
\end{tabular}
\end{table}

\begin{figure}[htbp]
\centering
\begin{tikzpicture}
\begin{axis}[
    ybar,
    bar width=16pt,
    width=0.88\textwidth,
    height=6.5cm,
    ymode=log,
    log basis y=10,
    ymin=10,
    ymax=200000,
    ylabel={Cycles (log scale)},
    symbolic x coords={LKL,kvm-v2 gadget,native,kvm-v2 exit,seccomp,ptrace},
    xtick=data,
    x tick label style={rotate=20,anchor=east,font=\small},
    nodes near coords,
    every node near coord/.append style={font=\tiny,above},
    enlarge x limits=0.18,
    grid=major,
    ymajorgrids=true,
    major grid style={draw=UMLGrid},
]
\addplot[fill=UMLBlue!60,draw=UMLBlue] coordinates {
    (LKL,18)
    (kvm-v2 gadget,90)
    (native,150)
    (kvm-v2 exit,36800)
    (seccomp,110000)
    (ptrace,9000)
};
\end{axis}
\end{tikzpicture}
\caption{Per-syscall round-trip cost comparison (log scale, cycles).
  The LSTAR gadget achieves $\sim$1,200$\times$ improvement over
  the \Seccomp{} backend for trivial syscalls.  LKL is shown for
  reference but sacrifices ring-split fidelity.}
\label{fig:syscall-cycles}
\end{figure}

\paragraph{The LSTAR gadget.}
The $\sim$25\,ns / $\sim$90-cycle result for \KVMvTwo{} with the
LSTAR gadget warrants explanation.  The gadget is an in-guest
\texttt{SYSCALL} handler installed at the LSTAR MSR address.  When
the guest executes \texttt{SYSCALL}, the CPU dispatches to the
gadget page \emph{without} a VMEXIT.  The gadget performs fast-path
handling (for trivial syscalls like \texttt{getpid}) and returns to
guest userspace via \texttt{SYSRET}.  Only syscalls requiring
kernel-side processing trigger a VMEXIT.  This achieves sub-native
latency because the native host syscall traverses speculative-execution
mitigations (KPTI, retpoline) that the in-guest gadget does not need.


% ------------------------------------------------------------------
\section{Cross-Host Benchmark Results}
\label{sec:bench-crosshost}
% ------------------------------------------------------------------

Table~\ref{tab:bench-crosshost} reports the complete cross-host
benchmark results across all three tiers.

\begin{table}[htbp]
\centering
\caption{Post-gadget \KVMvTwo{} speedup over seccomp across 8 hosts.}
\label{tab:bench-crosshost}
\begin{tabular}{@{}llrrr@{}}
\toprule
Host & CPU & Micro & Python & Stress \\
\midrule
s0 & Intel Core i9-12900K   & 284$\times$  & 3.73$\times$ & 7.26$\times$ \\
s1 & Intel Xeon E3-1225 v6  & 743$\times$  & 3.05$\times$ & 3.71$\times$ \\
s2 & Intel Xeon E3-1225 v5  & 813$\times$  & 2.99$\times$ & 3.94$\times$ \\
s3 & Intel Xeon W-2123      & 707$\times$  & 3.14$\times$ & 4.16$\times$ \\
s4 & Intel Core i5-12600K   & 193$\times$  & 2.99$\times$ & 8.72$\times$ \\
s5 & AMD Ryzen 7 7840HS     & 861$\times$  & 3.87$\times$ & 4.30$\times$ \\
s6 & AMD Ryzen 7 7840HS     & 908$\times$  & 4.06$\times$ & 4.32$\times$ \\
s7 & AMD Ryzen 7 7840HS     & 904$\times$  & 3.91$\times$ & 4.89$\times$ \\
\midrule
\multicolumn{2}{@{}l}{\textbf{Range}} & \textbf{193--908$\times$} & \textbf{2.99--4.06$\times$} & \textbf{3.71--8.72$\times$} \\
\bottomrule
\end{tabular}
\end{table}


\subsection{Micro Tier}

\begin{figure}[htbp]
\centering
\begin{tikzpicture}
\begin{axis}[
    ybar,
    bar width=14pt,
    width=0.88\textwidth,
    height=6cm,
    ylabel={Micro speedup over seccomp ($\times$)},
    symbolic x coords={s0,s1,s2,s3,s4,s5,s6,s7},
    xtick=data,
    ymin=0,
    nodes near coords,
    every node near coord/.append style={font=\small},
    enlarge x limits=0.12,
    grid=major,
    ymajorgrids=true,
    major grid style={draw=UMLGrid},
    xlabel={Host},
]
\addplot[fill=UMLTeal!70,draw=UMLTeal] coordinates {
    (s0,284) (s1,743) (s2,813) (s3,707) (s4,193)
    (s5,861) (s6,908) (s7,904)
};
\end{axis}
\end{tikzpicture}
\caption{Micro-tier (getpid-loop) speedup.  The variation across hosts
  reflects differences in seccomp notification overhead rather than
  KVM performance: the \KVMvTwo{} gadget cost is consistently 87--153
  cycles, while seccomp costs vary from $\sim$15,000 to $\sim$140,000
  cycles by host.}
\label{fig:bench-micro}
\end{figure}


\subsection{Application Tier (Python Startup)}

\begin{figure}[htbp]
\centering
\begin{tikzpicture}
\begin{axis}[
    ybar,
    bar width=14pt,
    width=0.88\textwidth,
    height=6cm,
    ylabel={Python startup speedup ($\times$)},
    symbolic x coords={s0,s1,s2,s3,s4,s5,s6,s7},
    xtick=data,
    ymin=0,
    ymax=5.5,
    nodes near coords,
    every node near coord/.append style={font=\small},
    enlarge x limits=0.12,
    grid=major,
    ymajorgrids=true,
    major grid style={draw=UMLGrid},
    xlabel={Host},
]
\addplot[fill=UMLGreen!70,draw=UMLGreen] coordinates {
    (s0,3.73) (s1,3.05) (s2,2.99) (s3,3.14) (s4,2.99)
    (s5,3.87) (s6,4.06) (s7,3.91)
};
\draw[dashed,UMLAmber,thick] (axis cs:s0,3.47) -- (axis cs:s7,3.47)
  node[pos=0.5,above,font=\footnotesize\color{UMLAmber}] {mean = 3.47$\times$};
\end{axis}
\end{tikzpicture}
\caption{\KVMvTwo{} Python startup speedup over seccomp.  The consistent
  3--4$\times$ band across diverse hardware confirms an architectural
  (backend mechanism) improvement, not a microarchitecture-specific one.}
\label{fig:bench-py-speedup}
\end{figure}


\subsection{Stress Tier}

\begin{figure}[htbp]
\centering
\begin{tikzpicture}
\begin{axis}[
    ybar,
    bar width=14pt,
    width=0.88\textwidth,
    height=6cm,
    ylabel={Stress speedup ($\times$)},
    symbolic x coords={s0,s1,s2,s3,s4,s5,s6,s7},
    xtick=data,
    ymin=0,
    ymax=10,
    nodes near coords,
    every node near coord/.append style={font=\small},
    enlarge x limits=0.12,
    grid=major,
    ymajorgrids=true,
    major grid style={draw=UMLGrid},
    xlabel={Host},
]
\addplot[fill=UMLAmber!70,draw=UMLAmber] coordinates {
    (s0,7.26) (s1,3.71) (s2,3.94) (s3,4.16) (s4,8.72)
    (s5,4.30) (s6,4.32) (s7,4.89)
};
\end{axis}
\end{tikzpicture}
\caption{Stress tier (mt-mini SMP T=8, ncpus=4) speedup.  The wider
  variance (3.71--8.72$\times$) reflects the mmap/page-fault-heavy
  workload where TLB shootdown cost and host scheduler behavior create
  per-host variation.  Zero strict memory failures on any host.}
\label{fig:bench-stress}
\end{figure}

The stress tier combines performance with correctness.  Each mt-mini
iteration allocates memory, writes patterns, forks or spawns threads,
and verifies memory content after concurrent modification.  A single
bit-flip constitutes a test failure.  The 3.71--8.72$\times$ speedup
comes from the broader vCPU and SMP design, not solely the gadget.


% ------------------------------------------------------------------
\section{Snapshot Performance}
\label{sec:bench-snapshot}
% ------------------------------------------------------------------

The \KVMvTwo{} snapshot benchmark exercises
\texttt{kvm\_v2\_snapshot\_alloc/capture/restore} with $N=64$
iterations of capture-once, restore-$N$-times.

\begin{table}[htbp]
\centering
\caption{Snapshot benchmark (kernel \texttt{9ccbf5e7bb58}, AMD Ryzen 7 7840HS, $N=64$).}
\label{tab:snapshot-perf}
\begin{tabular}{@{}lrrl@{}}
\toprule
Metric & Value & Design target & Margin \\
\midrule
Capture (one-shot)       & 35--39\,$\mu$s  & $<$\,50\,ms  & $\sim$1,300$\times$ under \\
Restore (median)         & 13\,$\mu$s      & $<$\,1\,ms   & $\sim$77$\times$ under \\
Restore (p95)            & 16--19\,$\mu$s  & ---           & --- \\
Restore (min)            & 12.5\,$\mu$s    & ---           & --- \\
Restore (max)            & 19--22\,$\mu$s  & ---           & --- \\
Mode                     & full            & ---           & memslot copy executed \\
\bottomrule
\end{tabular}
\end{table}

\begin{figure}[htbp]
\centering
\begin{tikzpicture}
\begin{axis}[
    ybar,
    bar width=20pt,
    width=0.72\textwidth,
    height=5.5cm,
    ylabel={Latency ($\mu$s)},
    symbolic x coords={capture,restore (med),restore (p95),restore (max)},
    xtick=data,
    x tick label style={font=\small},
    ymin=0,
    ymax=50,
    nodes near coords,
    every node near coord/.append style={font=\small},
    enlarge x limits=0.20,
    grid=major,
    ymajorgrids=true,
    major grid style={draw=UMLGrid},
]
\addplot[fill=UMLBlue!65,draw=UMLBlue] coordinates {
    (capture,38.9)
    (restore (med),13)
    (restore (p95),17.5)
    (restore (max),20.5)
};
\end{axis}
\end{tikzpicture}
\caption{Snapshot capture and restore latency.  The \SI{13}{\micro\second}
  median restore enables checkpoint-at-every-syscall workflows without
  budget pressure.}
\label{fig:bench-snapshot}
\end{figure}

\paragraph{Why so fast.}
The snapshot design avoids the three operations that dominate
checkpoint/restore cost in conventional systems: (1)~no per-restore
TLB shootdown (identity-mapped memslots); (2)~single host-VA-to-host-VA
\texttt{memcpy} for memslot restore (no page-table walk); (3)~fixed-size
\texttt{ioctl} chain for register/SREGS push (no variable-length data).

Record/replay KUnit suite (\texttt{kvm\_v2\_record}): 7/7 tests pass,
covering basic recording, state machine transitions, syscall observation,
strict replay, gadget bypass, RDTSC interception, and SIGALRM anchoring.


% ------------------------------------------------------------------
\section{CPython Substrate Parity}
\label{sec:bench-cpython}
% ------------------------------------------------------------------

The CPython parity gate runs 21 CPython 3.14 test modules on both
backends and requires identical results.

\begin{table}[htbp]
\centering
\caption{CPython substrate parity results.}
\label{tab:bench-cpython}
\begin{tabular}{@{}lcccc@{}}
\toprule
Backend & PASS & FAIL & XFAIL & Parity \\
\midrule
\KVMvTwo{}    & 25 & 3 & 3 & --- \\
Seccomp       & 25 & 3 & 3 & --- \\
\midrule
\textbf{Match} & \checkmark & \checkmark & \checkmark & \textbf{21/21} \\
\bottomrule
\end{tabular}
\end{table}

\noindent
The three FAIL results are known \UML{} limitations (not backend
regressions): tests depending on \texttt{/proc/self/maps} layout,
hardware timer resolution, or host-specific device nodes.  The three
XFAIL results are test-side expected failures unrelated to the backend.
The critical property is not that all tests pass, but that both
backends produce \emph{identical} results---confirming functional
equivalence of the \KVMvTwo{} and seccomp execution paths.

The wide parity gate extends this to 135 modules: 134/135 match
exactly, with a single \texttt{KVM\_BETTER} result (a test that
passes on \KVMvTwo{} but fails on seccomp due to timing sensitivity)
and zero regressions.


% ------------------------------------------------------------------
\section{Validation Gates}
\label{sec:bench-gates}
% ------------------------------------------------------------------

\begin{table}[htbp]
\centering
\caption{Key validation gates (May 2026, kernel HEAD on \texttt{next}).}
\label{tab:validation-gates}
\begin{tabular}{@{}lrrl@{}}
\toprule
Gate & Result & N & Notes \\
\midrule
cpython-parity (seccomp)          & 21/21 PARITY   & 21 modules  & UP and SMP \\
cpython-parity (kvm-v2)           & 21/21 PARITY   & 21 modules  & UP and SMP \\
Wide cpython-parity (kvm-v2)      & 134/135 PARITY & 135 modules & 1 KVM\_BETTER, 0 regression \\
Substrate gate (seccomp)          & 25/3/3         & 31 tests    & PASS/FAIL/XFAIL \\
Substrate gate (kvm-v2)           & 25/3/3         & 31 tests    & Matches seccomp exactly \\
mt-mini SMP (T=8, ncpus=4)        & 400/400        & 400 boots   & Wilson 95\% [99.0\%, 100.0\%] \\
threaded-fork-malloc               & 0/116,000      & 30 boots    & 0 CHILD\_FAIL across 116K forks \\
threaded-subprocess-wait           & 10/10          & 10 boots    & 0/4,000 subprocess failures \\
bench-micro getpid (kvm-v2)       & 89--105 cyc    & 8 hosts     & $\sim$1,000$\times$ vs seccomp \\
perf-py-startup (kvm-v2/seccomp)  & 1.10 ratio     & ---         & Gate ceiling: $\leq$1.20 \\
Phase J pilot soak                 & 240/240        & 3 workloads & memcheck/iocheck/stress-ng \\
Django Tier 3 (post-R14)           & 120/120        & loopback    & Wilson 95\% [96.90\%, 100.00\%] \\
\texttt{umlctl mission}           & 6/6 phases     & 100 soak    & 867s wall-clock; 0 panics \\
\bottomrule
\end{tabular}
\end{table}


% ------------------------------------------------------------------
\section{Profile Cost Model}
\label{sec:bench-profiles}
% ------------------------------------------------------------------

Table~\ref{tab:profile-cost} breaks down the per-syscall cost by
architectural layer for each profile, demonstrating the three-layer
composition property.

\begin{table}[htbp]
\centering
\caption{Per-syscall cost breakdown by profile and layer.}
\label{tab:profile-cost}
\begin{tabular}{@{}lrrrr@{}}
\toprule
Profile & Layer 1 (trap) & Layer 2 (gates) & Layer 3 (wraps) & Total \\
\midrule
prod-fast (KVM, no hooks)     & $\sim$70\,ns  & $\sim$1.5\,ns  & 0\,ns         & $\sim$72\,ns \\
prod-with-hooks (KVM, trace)  & $\sim$70\,ns  & $\sim$80\,ns   & 0\,ns         & $\sim$150\,ns \\
research (seccomp, full)      & $\sim$300\,ns & $\sim$80\,ns   & $\sim$100\,ns & $\sim$480\,ns \\
fuzz (seccomp, KCOV+KASAN)    & $\sim$300\,ns & $\sim$10\,ns   & $\sim$100\,ns & $\sim$410\,ns \\
sandbox (seccomp-only, none)  & $\sim$300\,ns & $\sim$1.5\,ns  & 0\,ns         & $\sim$302\,ns \\
\bottomrule
\end{tabular}
\end{table}

\begin{figure}[htbp]
\centering
\begin{tikzpicture}
\begin{axis}[
    ybar,
    bar width=9pt,
    width=0.92\textwidth,
    height=6.5cm,
    ymode=log,
    log basis y=10,
    ymin=4,
    ymax=3000,
    ylabel={Estimated getpid cost (ns, log scale)},
    symbolic x coords={library,native,prod-fast,prod-hooks,seccomp,sandbox,fuzz,research,fuzz-deep,time-travel,embedded},
    xtick=data,
    x tick label style={rotate=35,anchor=east,font=\footnotesize},
    nodes near coords,
    every node near coord/.append style={font=\tiny,rotate=90,anchor=west},
    enlarge x limits=0.08,
    grid=major,
    ymajorgrids=true,
    major grid style={draw=UMLGrid},
]
\addplot[fill=UMLBlue!70,draw=UMLBlue] coordinates {
    (library,6)
    (native,50)
    (prod-fast,80)
    (prod-hooks,85)
    (seccomp,400)
    (sandbox,505)
    (fuzz,550)
    (research,715)
    (fuzz-deep,1200)
    (time-travel,1200)
    (embedded,2000)
};
\end{axis}
\end{tikzpicture}
\caption{Profile cost model.  The value of the profile system is the
  ability to span from direct-call library mode (\SI{6}{\nano\second})
  through heavily instrumented research and time-travel builds
  (\SI{1.2}{\micro\second}) without forking the tree.  The
  \texttt{native} bar is a bare-metal host reference, not a UML
  profile.}
\label{fig:bench-profiles}
\end{figure}

\noindent
The key property: same source tree, radically different runtime
behaviors.  Layer~1 dominates (backend choice), Layer~2 is near-free
when off (static-key NOPs), Layer~3 is compile-time and orthogonal
to both.


% ------------------------------------------------------------------
\section{Summary of Performance Claims}
\label{sec:bench-summary}
% ------------------------------------------------------------------

\begin{table}[htbp]
\centering
\caption{Summary of validated performance claims.}
\label{tab:bench-summary}
\begin{tabularx}{\linewidth}{@{}p{0.26\linewidth}p{0.20\linewidth}p{0.22\linewidth}Y@{}}
\toprule
Claim & Target & Achieved & Methodology \\
\midrule
Gadget syscall latency     & $\sim$100\,ns         & $\sim$25\,ns / 90 cycles    & RDTSC-bracketed getpid \\
Python startup speedup     & faster than seccomp    & 2.99--4.06$\times$          & 8-host cross-comparison \\
Stress tier speedup        & faster than seccomp    & 3.71--8.72$\times$          & mt-mini SMP T=8 ncpus=4 \\
Micro tier speedup         & faster than seccomp    & 193--908$\times$            & getpid-loop \\
Snapshot capture           & $<$\,50\,ms            & 35--39\,$\mu$s              & kvm\_v2\_snapshot\_bench N=64 \\
Snapshot restore (median)  & $<$\,1\,ms             & 13\,$\mu$s                  & kvm\_v2\_snapshot\_bench N=64 \\
CPython parity             & identical results      & 21/21 match                 & cpython-parity gate \\
Soak pass rate             & 100\%                  & 100/100                     & umlctl mission \\
\bottomrule
\end{tabularx}
\end{table}

\paragraph{Interpretation.}
The performance story has four layers:
\begin{itemize}[leftmargin=2em]
  \item \textbf{Mechanism win.}  The micro tier proves that in-guest
    fast-path dispatch removes most seccomp signal/stub cost for
    trivial syscalls.
  \item \textbf{Workload win.}  The Python tier proves that real
    userland startup benefits from those trivial syscalls becoming
    cheap.
  \item \textbf{Concurrency win.}  The stress tier is mmap/page-fault
    heavy, so the gadget is less directly relevant; the continuing
    3.71--8.72$\times$ speedup comes from the broader vCPU and SMP
    design.
  \item \textbf{Snapshot win.}  The snapshot capture/restore numbers
    are within the sub-ms iteration budget by orders of magnitude.
    Record/replay sessions can now realistically checkpoint at every
    observed-syscall boundary without budget pressure.
\end{itemize}

\noindent
The correct benchmark unit is the ratio against seccomp on the same
host, not the raw absolute latency across machines.  That choice
removes most CPU-frequency and microarchitecture noise and keeps the
claim tied to a backend decision.
